
% Complex-field comparative replacement + dual memory representation

\subsection{Honest Address Resolution}

Resolve a permanent address only when both
\[
D_{(1)} < D_{\max}
\]
and
\[
D_{(2)}-D_{(1)} > \Delta_D.
\]
Otherwise return unresolved.

\subsection{Comparative Replacement}

For retained slot j:
\[
R_j =
U_j\sqrt{\log(1+n_j)}
\frac{1}{1+\alpha a_j}.
\]

For candidate c:
\[
C_c =
U_c(0.65+0.35S_c)(0.75+0.25Q_c).
\]

If the bank is full, replace the weakest retained slot only when
\[
C_c-\min_j R_j > m_R.
\]

\subsection{Dual Representation}

Maintain separately
\[
\bar{\mathbf v}_c
=
\text{averaged robust address vector}
\]
and
\[
\chi_c
=
\text{best topology-qualified complex observation}.
\]

Do not arithmetic-average complex phase templates merely because their phase-agnostic
address vectors belong to the same candidate.

A candidate may be ADDRESS_READY while remaining TEMPLATE_UNRESOLVED.

\subsection{Robust Winding Validation}

For a finite loop C enclosing a recovered core,
\[
W_C =
\frac{1}{2\pi}
\sum_{k\in C}
\arg(\psi_{k+1}\psi_k^*).
\]

Use finite-loop winding for recall validation rather than exact global microscopic
plaquette-defect equality near amplitude singularities.
